% ===== PRÉAMBULE CANONIQUE — à coller VERBATIM en tête de CHAQUE .tex =====
\documentclass[11pt,a4paper]{article}
\usepackage[utf8]{inputenc}\usepackage[T1]{fontenc}\usepackage{lmodern}
\usepackage[french]{babel}
\usepackage{amsmath,amssymb,mathtools}\usepackage{icomma}
\usepackage[version=4]{mhchem}          % \ce{...} : équations & formules
\usepackage{siunitx}\sisetup{locale=FR} % \qty{}{}, \si{}, \ang{}
\usepackage{chemfig}                    % \chemfig{...} : structures développées
\usepackage{xcolor}\usepackage{colortbl}
\usepackage{tikz}\usepackage{pgfplots}\pgfplotsset{compat=1.18}
\usepackage[most]{tcolorbox}\usepackage{enumitem}\usepackage{array,booktabs}
\usepackage[a4paper,margin=2.2cm]{geometry}
\usetikzlibrary{arrows.meta,calc,angles,quotes,positioning,patterns,backgrounds,shapes.geometric}
\definecolor{primary}{RGB}{222,49,99}\definecolor{primarydark}{RGB}{150,22,55}
\definecolor{defcol}{RGB}{0,114,178}\definecolor{methcol}{RGB}{52,92,148}
\definecolor{excol}{RGB}{0,135,90}\definecolor{attncol}{RGB}{200,40,55}
\tcbset{boxbase/.style={enhanced,breakable,boxrule=0.9pt,arc=2.5pt,
  left=3mm,right=3mm,top=2.3mm,bottom=2.3mm,fonttitle=\bfseries,coltitle=white}}
% --- boîtes TUTORIEL (noms EXACTS, ne pas renommer) ---
\newtcolorbox{cle}[1][]{boxbase,colback=defcol!6,colframe=defcol,
  title={\textsc{À retenir}\ifx\relax#1\relax\relax\else\textmd{ : \itshape #1}\fi}}
\newtcolorbox{methode}[1][]{boxbase,colback=methcol!7,colframe=methcol,sharp corners,
  title={\textsc{Méthode}\ifx\relax#1\relax\relax\else\textmd{ --- \itshape #1}\fi}}
\newtcolorbox{exemplebox}[1][]{boxbase,colback=excol!5,colframe=excol,
  title={\textsc{Exemple}\ifx\relax#1\relax\relax\else\textmd{ : \itshape #1}\fi}}
\newtcolorbox{piege}[1][]{boxbase,colback=attncol!6,colframe=attncol,
  title={\textsc{Piège}\ifx\relax#1\relax\relax\else\textmd{ : \itshape #1}\fi}}
\newtcolorbox{remarque}[1][]{boxbase,colback=black!4,colframe=black!45,
  title={\textsc{Remarque}\ifx\relax#1\relax\relax\else\textmd{ : \itshape #1}\fi}}
% --- boîtes EXERCICE / CORRIGÉ (noms EXACTS, auto-comptées) ---
\newtcolorbox[auto counter]{exo}[1][]{boxbase,colback=primary!4,colframe=primary,
  title={\textsc{Exercice}~\thetcbcounter\ifx\relax#1\relax\relax\else\textmd{ --- \itshape #1}\fi}}
\newtcolorbox[auto counter]{corrige}[1][]{boxbase,colback=excol!4,colframe=excol,
  title={\textsc{Corrigé}~\thetcbcounter\ifx\relax#1\relax\relax\else\textmd{ --- \itshape #1}\fi}}
\newcommand{\R}{\mathbb{R}}\newcommand{\N}{\mathbb{N}}\newcommand{\Z}{\mathbb{Z}}\newcommand{\ds}{\displaystyle}
\begin{document}
% THEME_TITLE: Reconnaître les groupements fonctionnels

Reconnaître les fonctions d'une molécule, c'est déjà connaître sa chimie. C'est \emph{la} compétence testée au concours : on donne une structure, on demande quelles fonctions elle porte. Toute la méthode tient en deux réflexes : chercher les \textbf{carbonyles} $\ce{C=O}$, puis les \textbf{hétéroatomes isolés}.

\begin{cle}[Un carbonyle $\ce{C=O}$, cinq fonctions selon le voisinage]
\begin{center}\small
\renewcommand{\arraystretch}{1.25}
\begin{tabular}{lll}
\toprule
voisinage du $\ce{C=O}$ & fonction & suffixe \\
\midrule
$-\ce{C(=O)}-\ce{H}$ (un H) & \textbf{aldéhyde} & \textit{-al} \\
$-\ce{C(=O)}-$ (deux C) & \textbf{cétone} & \textit{-one} \\
$-\ce{C(=O)}-\ce{OH}$ & \textbf{acide carboxylique} & \textit{acide \dots-oïque} \\
$-\ce{C(=O)}-\ce{O}-\ce{R}$ & \textbf{ester} & \textit{-oate d'alkyle} \\
$-\ce{C(=O)}-\ce{N}$ & \textbf{amide} & \textit{-amide} \\
\bottomrule
\end{tabular}
\end{center}
\end{cle}

\begin{cle}[Les fonctions à hétéroatome « isolé » (pas de carbonyle)]
\begin{itemize}[leftmargin=1.6em,itemsep=0.5pt]
  \item \textbf{alcool} $-\ce{OH}$ sur un carbone saturé ; \quad \textbf{éther} $-\ce{O}-$ pontant deux carbones ;
  \item \textbf{amine} $-\ce{N}$ liée à des carbones \emph{sans} carbonyl voisin ; \quad \textbf{thiol} $-\ce{SH}$ ;
  \item \textbf{nitrile} $-\ce{C#N}$ ; \quad \textbf{alcène} $\ce{C=C}$ ; \quad \textbf{alcyne} $\ce{C#C}$.
\end{itemize}
\end{cle}

\begin{methode}[Lire une structure, fonction par fonction]
\begin{enumerate}[leftmargin=1.6em,itemsep=0.5pt]
  \item Repérer chaque $\ce{C=O}$, puis regarder ce qui l'entoure : H / C+C / OH / O--R / N $\rightarrow$ aldéhyde / cétone / acide / ester / amide.
  \item Repérer les hétéroatomes \emph{restants} : $-\ce{OH}$ (alcool), $-\ce{O}-$ pont (éther), $-\ce{N}$ sans carbonyle (amine), $-\ce{SH}$ (thiol).
  \item Repérer les liaisons multiples C--C ($\ce{C=C}$ alcène, $\ce{C#C}$ alcyne).
  \item Pour alcools et amines, préciser la classe (primaire / secondaire / tertiaire) d'après le carbone porteur.
\end{enumerate}
\end{methode}

\begin{piege}[Les trois confusions qui coûtent des points]
\begin{itemize}[leftmargin=1.6em,itemsep=0.5pt]
  \item \textbf{amine ou amide ?} Un azote collé à un $\ce{C=O}$ ($-\ce{CO-N}$) est une \textbf{amide}, jamais une amine. L'amine exige un azote \emph{sans} carbonyle voisin.
  \item \textbf{ester ou éther ?} Dans $-\ce{CO-O-R}$ l'oxygène pontant fait partie de l'\textbf{ester} : ce n'est \emph{pas} une fonction éther (piège classique de l'aspirine).
  \item \textbf{alcool ou éther ?} $-\ce{OH}$ terminal $=$ alcool ; $-\ce{O}-$ entre deux carbones $=$ éther. Mêmes atomes ($\ce{C2H6O}$), fonctions différentes (isomères de fonction).
\end{itemize}
\end{piege}

\end{document}
